\documentclass[11pt]{article}
% The preamble every project's paper shares: % The preamble every project's paper shares: % The preamble every project's paper shares: \input{../../preamble} from
% projects/<project>/paper/. Boards come from chessboard and xskak, plots from
% pgfplots; every number a paper prints comes from its generated.tex.
\usepackage[LSBC4,T1]{fontenc}
\usepackage[utf8]{inputenc}
\usepackage{lmodern}
\usepackage[margin=1in]{geometry}
\usepackage{microtype}
\usepackage{amsmath}
\usepackage{booktabs}
\usepackage{tabularx}
\usepackage{array}
\usepackage{xcolor}
\usepackage[obeyspaces,spaces]{url}
\usepackage{xskak}
\usepackage{pgfplots}
\pgfplotsset{compat=1.18}
\usepackage{caption}
\usepackage{subcaption}
\usepackage{enumitem}
\usepackage[section]{placeins}
\usepackage[hidelinks]{hyperref}

\setchessboard{boardfontencoding=LSBC4, boardfontfamily=skaknew, blackfieldcolor=black!22,
  whitefieldcolor=white, setfontcolors, showmover=false, label=true, labelfontsize=6pt,
  boardfontsize=14pt}

% A move or line in UCI, and a position in FEN, both breakable at a slash.
\newcommand{\uci}[1]{\texttt{#1}}
\newcommand{\fen}[1]{{\footnotesize\path{#1}}}
\newcolumntype{L}{>{\raggedright\arraybackslash}X}

\setlength{\parskip}{0.4em}
\setlength{\parindent}{0pt}
 from
% projects/<project>/paper/. Boards come from chessboard and xskak, plots from
% pgfplots; every number a paper prints comes from its generated.tex.
\usepackage[LSBC4,T1]{fontenc}
\usepackage[utf8]{inputenc}
\usepackage{lmodern}
\usepackage[margin=1in]{geometry}
\usepackage{microtype}
\usepackage{amsmath}
\usepackage{booktabs}
\usepackage{tabularx}
\usepackage{array}
\usepackage{xcolor}
\usepackage[obeyspaces,spaces]{url}
\usepackage{xskak}
\usepackage{pgfplots}
\pgfplotsset{compat=1.18}
\usepackage{caption}
\usepackage{subcaption}
\usepackage{enumitem}
\usepackage[section]{placeins}
\usepackage[hidelinks]{hyperref}

\setchessboard{boardfontencoding=LSBC4, boardfontfamily=skaknew, blackfieldcolor=black!22,
  whitefieldcolor=white, setfontcolors, showmover=false, label=true, labelfontsize=6pt,
  boardfontsize=14pt}

% A move or line in UCI, and a position in FEN, both breakable at a slash.
\newcommand{\uci}[1]{\texttt{#1}}
\newcommand{\fen}[1]{{\footnotesize\path{#1}}}
\newcolumntype{L}{>{\raggedright\arraybackslash}X}

\setlength{\parskip}{0.4em}
\setlength{\parindent}{0pt}
 from
% projects/<project>/paper/. Boards come from chessboard and xskak, plots from
% pgfplots; every number a paper prints comes from its generated.tex.
\usepackage[LSBC4,T1]{fontenc}
\usepackage[utf8]{inputenc}
\usepackage{lmodern}
\usepackage[margin=1in]{geometry}
\usepackage{microtype}
\usepackage{amsmath}
\usepackage{booktabs}
\usepackage{tabularx}
\usepackage{array}
\usepackage{xcolor}
\usepackage[obeyspaces,spaces]{url}
\usepackage{xskak}
\usepackage{pgfplots}
\pgfplotsset{compat=1.18}
\usepackage{caption}
\usepackage{subcaption}
\usepackage{enumitem}
\usepackage[section]{placeins}
\usepackage[hidelinks]{hyperref}

\setchessboard{boardfontencoding=LSBC4, boardfontfamily=skaknew, blackfieldcolor=black!22,
  whitefieldcolor=white, setfontcolors, showmover=false, label=true, labelfontsize=6pt,
  boardfontsize=14pt}

% A move or line in UCI, and a position in FEN, both breakable at a slash.
\newcommand{\uci}[1]{\texttt{#1}}
\newcommand{\fen}[1]{{\footnotesize\path{#1}}}
\newcolumntype{L}{>{\raggedright\arraybackslash}X}

\setlength{\parskip}{0.4em}
\setlength{\parindent}{0pt}

% Written by projects/break-ready/paper/generate.mjs; do not edit by hand.

\newcommand{\AppCommit}{f4bd63f19b}
\newcommand{\GradeNodes}{250,000}
\newcommand{\Window}{3}
\newcommand{\Draws}{2,000}
\newcommand{\Seed}{20261020}
\newcommand{\AimedMin}{two}
\newcommand{\MinArm}{60}
\newcommand{\MinCell}{30}
\newcommand{\PositionsFirst}{491}
\newcommand{\GamesFirst}{114}
\newcommand{\BreaksFirst}{716}
\newcommand{\PositionsSecond}{418}
\newcommand{\GamesSecond}{103}
\newcommand{\BreaksSecond}{553}
\newcommand{\StrongPositions}{1,200}
\newcommand{\StrongBreaks}{1,417}
\newcommand{\BreaksAll}{1,269}
\newcommand{\PositionsAll}{909}
\newcommand{\Floor}{2200}
\newcommand{\CoverageTable}{%
\begin{table}[ht]
\centering\small
\setlength{\tabcolsep}{4pt}
\begin{tabular}{@{}llll@{}}
\toprule
& His, first half & His, second half & Strong players \\
\midrule
Breaks (positions) & 716 (491) & 553 (418) & 1,417 (1,200) \\
\midrule
1. Pieces aimed & 359 {\scriptsize(50.1\%)} & 294 {\scriptsize(53.2\%)} & 507 {\scriptsize(35.8\%)} \\
2. Rooks ready & 567 {\scriptsize(79.2\%)} & 474 {\scriptsize(85.7\%)} & 1,171 {\scriptsize(82.6\%)} \\
3. A passer or an exposed king & 50 {\scriptsize(7.0\%)} & 42 {\scriptsize(7.6\%)} & 165 {\scriptsize(11.6\%)} \\
4. Force on that side & 369 {\scriptsize(51.5\%)} & 335 {\scriptsize(60.6\%)} & 787 {\scriptsize(55.5\%)} \\
5. The blunder check: safe & 446 {\scriptsize(62.3\%)} & 344 {\scriptsize(62.2\%)} & 758 {\scriptsize(53.5\%)} \\
\midrule
All five & 4 {\scriptsize(0.6\%)} & 4 {\scriptsize(0.7\%)} & 11 {\scriptsize(0.8\%)} \\
Safe, with the pieces ready & 157 of 446 & 136 of 344 & 145 of 758 \\
\bottomrule
\end{tabular}
\caption{Counted before the rule was fixed, from the board alone, with no score read: the breaks each question passes, all five, and the safe breaks with the pieces ready (questions 1, 2 and 4), among the safe breaks.}
\label{tab:coverage}
\end{table}
}
\newcommand{\AllFive}{eight}
\newcommand{\AllFiveStrong}{eleven}
\newcommand{\PurposeShare}{7}
\newcommand{\ReadyFirst}{157}
\newcommand{\SafeFirst}{446}
\newcommand{\ReadySecond}{136}
\newcommand{\SafeSecond}{344}
\newcommand{\ReadyStrong}{145}
\newcommand{\SafeStrong}{758}
\newcommand{\Verdict}{no signal}
\newcommand{\LiftFirst}{3.0}
\newcommand{\LiftFirstCi}{$-$5.7 to 11.2}
\newcommand{\CheckFirst}{12.5}
\newcommand{\CheckFirstCi}{2.5 to 22.1}
\newcommand{\NullFirst}{$-$10.7}
\newcommand{\NullFirstCi}{$-$18.2 to $-$2.5}
\newcommand{\ControlFirst}{36.3}
\newcommand{\ControlFirstCi}{30.8 to 41.5}
\newcommand{\ReadySoundFirst}{26.1}
\newcommand{\RushedSoundFirst}{26.0}
\newcommand{\SafeSoundFirst}{26.0}
\newcommand{\UnsafeSoundFirst}{13.7}
\newcommand{\LiftSecond}{2.3}
\newcommand{\LiftSecondCi}{$-$6.0 to 11.8}
\newcommand{\CheckSecond}{15.9}
\newcommand{\CheckSecondCi}{6.1 to 26.0}
\newcommand{\NullSecond}{9.1}
\newcommand{\NullSecondCi}{0.2 to 18.9}
\newcommand{\ControlSecond}{37.6}
\newcommand{\ControlSecondCi}{31.7 to 43.5}
\newcommand{\ReadySoundSecond}{33.8}
\newcommand{\RushedSoundSecond}{29.8}
\newcommand{\SafeSoundSecond}{31.4}
\newcommand{\UnsafeSoundSecond}{12.9}
\newcommand{\StrongLift}{$-$7.0}
\newcommand{\StrongLiftCi}{$-$15.8 to 1.6}
\newcommand{\StrongCheck}{20.8}
\newcommand{\StrongCheckCi}{16.4 to 25.1}
\newcommand{\StrongControl}{51.8}
\newcommand{\StrongControlCi}{48.7 to 54.8}
\newcommand{\StrongReadySound}{29.0}
\newcommand{\StrongRushedSound}{33.3}
\newcommand{\StrongSafeSound}{32.5}
\newcommand{\StrongUnsafeSound}{12.3}
\newcommand{\ControlTable}{%
\begin{table}[ht]
\centering\small
\setlength{\tabcolsep}{4pt}
\begin{tabular}{@{}llll@{}}
\toprule
& His, first half & His, second half & Strong players \\
\midrule
1. The move played minus a random move & 36.3 {\scriptsize(30.8, 41.5)} & 37.6 {\scriptsize(31.7, 43.5)} & 51.8 {\scriptsize(48.7, 54.8)} \\
2. The blunder check's lift & 12.5 {\scriptsize(2.5, 22.1)} & 15.9 {\scriptsize(6.1, 26.0)} & 20.8 {\scriptsize(16.4, 25.1)} \\
3. A random reading's lift & $-$10.7 {\scriptsize($-$18.2, $-$2.5)} & 9.1 {\scriptsize(0.2, 18.9)} & -- \\
\midrule
Pieces ready, the lift & 3.0 {\scriptsize($-$5.7, 11.2)} & 2.3 {\scriptsize($-$6.0, 11.8)} & $-$7.0 {\scriptsize($-$15.8, 1.6)} \\
Safe breaks: ready / rushed & 157 / 289 & 136 / 208 & 145 / 613 \\
Sound: ready / rushed (\%) & 26.1 / 26.0 & 33.8 / 29.8 & 29.0 / 33.3 \\
Sound: safe / not safe (\%) & 26.0 / 13.7 & 31.4 / 12.9 & 32.5 / 12.3 \\
\bottomrule
\end{tabular}
\caption{The controls and the rule's lift, in points, with 95\% intervals resampling his games, or the strong players' positions. The strong players' column is reported, not ruled on. Below, the share of breaks that are sound on each side, in \%, before the groups are weighed.}
\label{tab:controls}
\end{table}
}
\newcommand{\QuestionTable}{%
\begin{table}[ht]
\centering\footnotesize
\setlength{\tabcolsep}{3pt}
\begin{tabular}{@{}lrllll@{}}
\toprule
Question & \begin{tabular}[b]{@{}l@{}}His breaks: \\ pass / fail\end{tabular} & \begin{tabular}[b]{@{}l@{}}His, halves \\ together\end{tabular} & \begin{tabular}[b]{@{}l@{}}His, first \\ half\end{tabular} & \begin{tabular}[b]{@{}l@{}}His, second \\ half\end{tabular} & \begin{tabular}[b]{@{}l@{}}Strong \\ players\end{tabular} \\
\midrule
Pieces aimed & 653 / 616 & 6.8 {\scriptsize(1.7, 12.6)} & 8.9 {\scriptsize(2.4, 15.9)} & 4.0 {\scriptsize($-$3.9, 13.2)} & $-$0.8 {\scriptsize($-$5.4, 3.7)} \\
Rooks ready & 1,041 / 228 & 3.6 {\scriptsize($-$3.5, 9.8)} & 4.8 {\scriptsize($-$4.7, 12.6)} & 1.0 {\scriptsize($-$9.5, 10.8)} & 0.5 {\scriptsize($-$5.4, 6.3)} \\
A passer or an exposed king & 92 / 1,177 & 9.1 {\scriptsize($-$7.3, 27.1)} & 24.0 {\scriptsize(4.4, 52.1)} & $-$4.7 {\scriptsize($-$25.8, 22.6)} & 1.5 {\scriptsize($-$6.7, 9.6)} \\
Force on that side & 704 / 565 & 9.2 {\scriptsize(3.5, 14.4)} & 8.2 {\scriptsize(1.5, 15.3)} & 9.2 {\scriptsize(0.3, 18.2)} & 1.8 {\scriptsize($-$2.3, 6.1)} \\
The blunder check: safe & 790 / 479 & 14.2 {\scriptsize(7.7, 20.4)} & 12.5 {\scriptsize(3.0, 22.2)} & 15.9 {\scriptsize(6.4, 26.0)} & 20.8 {\scriptsize(16.3, 25.2)} \\
The four but purpose, together & 293 / 976 & 9.4 {\scriptsize(3.9, 15.1)} & 7.3 {\scriptsize($-$0.1, 15.0)} & 10.4 {\scriptsize(2.5, 18.2)} & 6.1 {\scriptsize($-$1.6, 14.1)} \\
\bottomrule
\end{tabular}
\caption{Each question alone over all the breaks, safe or not, and the four questions other than purpose together: the lift, in points, with 95\% intervals.}
\label{tab:questions}
\end{table}
}
\newcommand{\QAimed}{6.8}
\newcommand{\QAimedCi}{1.7 to 12.6}
\newcommand{\QAimedStrong}{$-$0.8}
\newcommand{\QRooks}{3.6}
\newcommand{\QRooksCi}{$-$3.5 to 9.8}
\newcommand{\QRooksStrong}{0.5}
\newcommand{\QPurpose}{9.1}
\newcommand{\QPurposeCi}{$-$7.3 to 27.1}
\newcommand{\QPurposeStrong}{1.5}
\newcommand{\QForce}{9.2}
\newcommand{\QForceCi}{3.5 to 14.4}
\newcommand{\QForceStrong}{1.8}
\newcommand{\QSafe}{14.2}
\newcommand{\QSafeCi}{7.7 to 20.4}
\newcommand{\QSafeStrong}{20.8}
\newcommand{\QAllButPurpose}{9.4}
\newcommand{\QAllButPurposeCi}{3.9 to 15.1}
\newcommand{\QAllButPurposeStrong}{6.1}
\newcommand{\CountTable}{%
\begin{table}[ht]
\centering\small
\setlength{\tabcolsep}{4pt}
\begin{tabular}{@{}lrl@{\hspace{1.5em}}lrlrl@{}}
\toprule
\multicolumn{3}{@{}l}{All his breaks, by questions passed} & \multicolumn{5}{l}{The safe breaks, by piece questions passed} \\
\cmidrule(r{1.5em}){1-3}\cmidrule(l){4-8}
Of five & Breaks & Sound (\%) & Of three & \begin{tabular}[b]{@{}l@{}}His \\ breaks\end{tabular} & His, sound (\%) & \begin{tabular}[b]{@{}l@{}}Strong \\ breaks\end{tabular} & Strong, sound (\%) \\
\midrule
0 & 45 & 0.0 {\scriptsize(0.0, 0.0)} & 0 & 37 & 13.5 {\scriptsize(3.4, 23.1)} & 57 & 26.3 {\scriptsize(15.8, 37.9)} \\
1 & 189 & 14.3 {\scriptsize(8.6, 20.8)} & 1 & 185 & 20.5 {\scriptsize(14.3, 27.4)} & 205 & 34.1 {\scriptsize(27.5, 40.6)} \\
2 & 345 & 18.0 {\scriptsize(14.0, 22.6)} & 2 & 275 & 34.2 {\scriptsize(27.4, 41.3)} & 351 & 33.9 {\scriptsize(29.2, 39.0)} \\
3 & 367 & 25.9 {\scriptsize(20.7, 31.0)} & 3 & 293 & 29.7 {\scriptsize(23.8, 35.8)} & 145 & 29.0 {\scriptsize(21.8, 36.3)} \\
4 & 315 & 32.1 {\scriptsize(25.9, 38.0)} & & & & & \\
5 & 8 & \emph{n=8} & & & & & \\
\bottomrule
\end{tabular}
\caption{The share of breaks that are sound, in \%, with 95\% intervals. Left: all his breaks, both halves together, by how many of the five questions they pass (fewer than 30 breaks shows only their number). Right, added after the run: the safe breaks, his and the strong players', by how many of the three piece questions they pass.}
\label{tab:count}
\end{table}
}
\newcommand{\CountNone}{0.0}
\newcommand{\CountOne}{14.3}
\newcommand{\CountTwo}{18.0}
\newcommand{\CountThree}{25.9}
\newcommand{\CountFour}{32.1}
\newcommand{\GroupTable}{%
\begin{table}[ht]
\centering\small
\setlength{\tabcolsep}{4pt}
\begin{tabular}{@{}lrlll@{}}
\toprule
Break & Ready / rushed & Ready, sound (\%) & Rushed, sound (\%) & Lift (points) \\
\midrule
London pyramid, c3-c4 & 113 / 183 & 30.1 {\scriptsize(21.2, 38.9)} & 28.4 {\scriptsize(20.7, 37.1)} & 1.7 {\scriptsize($-$8.3, 10.8)} \\
London pyramid, e3-e4 & 93 / 107 & 38.7 {\scriptsize(28.0, 50.0)} & 30.8 {\scriptsize(21.2, 41.6)} & 7.9 {\scriptsize($-$6.9, 21.5)} \\
Advance chain, \ldots{}f7-f6 & 37 / 67 & 13.5 {\scriptsize(3.0, 29.4)} & 9.0 {\scriptsize(2.1, 18.6)} & 4.6 {\scriptsize($-$6.6, 19.5)} \\
\bottomrule
\end{tabular}
\caption{His safe breaks, both halves together, for each group with at least 30 breaks with the pieces ready and 30 rushed: the share sound on each side and the difference, with 95\% intervals.}
\label{tab:groups}
\end{table}
}
\newcommand{\GroupsOther}{19}
\newcommand{\PlayedReady}{7.2}
\newcommand{\PlayedRushed}{8.7}
\newcommand{\PlayedUnsafe}{3.1}
\newcommand{\StrongPlayedReady}{16.6}
\newcommand{\StrongPlayedRushed}{14.0}
\newcommand{\StrongPlayedUnsafe}{4.2}
\newcommand{\ChoiceLift}{$-$1.0}
\newcommand{\ChoiceLiftCi}{$-$4.6 to 2.9}
\newcommand{\AgainstFirst}{138}
\newcommand{\AgainstReadyFirst}{26.1}
\newcommand{\AgainstOwnFirst}{54.3}
\newcommand{\AgainstGainFirst}{$-$28.3}
\newcommand{\AgainstGainFirstCi}{$-$37.7 to $-$18.6}
\newcommand{\AgainstSecond}{111}
\newcommand{\AgainstReadySecond}{33.3}
\newcommand{\AgainstOwnSecond}{57.7}
\newcommand{\AgainstGainSecond}{$-$24.3}
\newcommand{\AgainstGainSecondCi}{$-$32.6 to $-$16.2}
\newcommand{\ChoiceTable}{%
\begin{table}[ht]
\centering\small
\setlength{\tabcolsep}{4pt}
\begin{tabular}{@{}llll@{}}
\toprule
& Safe, pieces ready & Safe, rushed & Not safe \\
\midrule
He played the break (\%) & 7.2 & 8.7 & 3.1 \\
Strong players played it (\%) & 16.6 & 14.0 & 4.2 \\
\midrule
& Positions & \begin{tabular}[b]{@{}l@{}}The ready \\ breaks, sound\end{tabular} & \begin{tabular}[b]{@{}l@{}}His move, \\ sound\end{tabular} \\
\midrule
His, first half & 138 & 26.1 {\scriptsize(18.1, 34.3)} & 54.3 {\scriptsize(45.8, 62.6)} \\
His, second half & 111 & 33.3 {\scriptsize(25.2, 42.3)} & 57.7 {\scriptsize(49.5, 67.0)} \\
\bottomrule
\end{tabular}
\caption{Above: how often the break was the move played, by what the questions say of it. Below: his positions with a safe break whose pieces are ready, the share of those breaks that are sound and of his moves there, in \%, with 95\% intervals.}
\label{tab:choice}
\end{table}
}
\newcommand{\CalSpreadFirst}{4.4}
\newcommand{\CalBootFirst}{4.3}
\newcommand{\CalExcludesFirst}{three}
\newcommand{\CalPFirst}{0.54}
\newcommand{\CalSpreadSecond}{5.2}
\newcommand{\CalBootSecond}{5.2}
\newcommand{\CalExcludesSecond}{six}
\newcommand{\CalPSecond}{0.68}
\newcommand{\CalReadings}{100}
\newcommand{\CalPermutations}{1,000}
\newcommand{\NullSdFirst}{2.4}
\newcommand{\NullSdSecond}{1.7}
\newcommand{\SafeQuestionTable}{%
\begin{table}[ht]
\centering\footnotesize
\setlength{\tabcolsep}{3pt}
\begin{tabular}{@{}lrllll@{}}
\toprule
Question & \begin{tabular}[b]{@{}l@{}}His safe breaks: \\ pass / fail\end{tabular} & \begin{tabular}[b]{@{}l@{}}His, halves \\ together\end{tabular} & \begin{tabular}[b]{@{}l@{}}His, first \\ half\end{tabular} & \begin{tabular}[b]{@{}l@{}}His, second \\ half\end{tabular} & \begin{tabular}[b]{@{}l@{}}Strong \\ players\end{tabular} \\
\midrule
Pieces aimed & 465 / 325 & 7.4 {\scriptsize(1.1, 14.8)} & 9.9 {\scriptsize(1.6, 18.3)} & 4.0 {\scriptsize($-$6.3, 17.7)} & $-$6.0 {\scriptsize($-$13.5, 1.2)} \\
Rooks ready & 674 / 116 & $-$0.9 {\scriptsize($-$11.0, 8.4)} & $-$3.6 {\scriptsize($-$17.8, 7.9)} & 1.2 {\scriptsize($-$11.7, 16.2)} & $-$1.4 {\scriptsize($-$9.8, 6.9)} \\
Force on that side & 475 / 315 & 8.4 {\scriptsize(0.6, 15.8)} & 8.9 {\scriptsize($-$1.2, 17.8)} & 5.5 {\scriptsize($-$7.1, 17.8)} & 3.2 {\scriptsize($-$3.9, 10.2)} \\
\midrule
\multicolumn{6}{@{}l}{The intervals, checked on random readings drawn as the null draws its own} \\
& \begin{tabular}[b]{@{}l@{}}Spread of \\ random lifts\end{tabular} & \begin{tabular}[b]{@{}l@{}}The bootstrap's \\ spread\end{tabular} & \begin{tabular}[b]{@{}l@{}}Intervals \\ leaving out 0\end{tabular} & \begin{tabular}[b]{@{}l@{}}Pieces ready's \\ lift\end{tabular} & \begin{tabular}[b]{@{}l@{}}Random lifts \\ as far from 0\end{tabular} \\
\midrule
His, first half & 4.4 & 4.3 & 3 of 100 & 3.0 & 0.54 \\
His, second half & 5.2 & 5.2 & 6 of 100 & 2.3 & 0.68 \\
\bottomrule
\end{tabular}
\caption{Added after the run. Above: each piece question alone among the safe breaks, the lift in points with 95\% intervals. Below: the spread of the lift, in points, over 1,000 random readings that pass as many safe breaks in each group as pieces ready does, the bootstrap's spread for 100 of them, how many of those 100 intervals leave out 0, and the share of random readings whose lift is at least as far from 0 as pieces ready's.}
\label{tab:safequestions}
\end{table}
}
\newcommand{\SAimed}{7.4}
\newcommand{\SAimedCi}{1.1 to 14.8}
\newcommand{\SAimedStrong}{$-$6.0}
\newcommand{\SRooks}{$-$0.9}
\newcommand{\SRooksCi}{$-$11.0 to 8.4}
\newcommand{\SRooksStrong}{$-$1.4}
\newcommand{\SForce}{8.4}
\newcommand{\SForceCi}{0.6 to 15.8}
\newcommand{\SForceStrong}{3.2}
\newcommand{\SafeTwo}{34.2}
\newcommand{\SafeThree}{29.7}
\newcommand{\EaFen}{\fen{1k1r3r/pppq1pbp/5np1/3p4/1n1P1B2/1N1BPQ1P/PP3PP1/R4RK1 w - - 2 14}}
\newcommand{\EaMover}{White}
\newcommand{\EaPhase}{middlegame}
\newcommand{\EaStructure}{the Carlsbad}
\newcommand{\EaBreak}{14.~e4}
\newcommand{\EaBreakSound}{not sound}
\newcommand{\EaWpBreak}{36}
\newcommand{\EaBest}{14.~\symbishop{}e2}
\newcommand{\EaWpBest}{69}
\newcommand{\EaPlayed}{14.~\symbishop{}e2}
\newcommand{\EaWpPlayed}{68}
\newcommand{\EaPlayedSound}{sound}
\newcommand{\EaAimed}{\symbishop{}d3 and \symqueen{}f3}
\newcommand{\EaAimedOk}{yes}
\newcommand{\EaUndeveloped}{}
\newcommand{\EaFiles}{d-file and e-file}
\newcommand{\EaRooks}{\symrook{}a1 and \symrook{}f1}
\newcommand{\EaRooksOk}{yes}
\newcommand{\EaForceUs}{20}
\newcommand{\EaForceThem}{17}
\newcommand{\EaForceOk}{yes}
\newcommand{\EaBoard}{%
\chessboard[setfen={1k1r3r/pppq1pbp/5np1/3p4/1n1P1B2/1N1BPQ1P/PP3PP1/R4RK1 w - - 2 14}, boardfontsize=14pt, pgfstyle=straightmove, arrow=to, linewidth=0.07em, color=blue!70!black, markmoves={d3-e2}, pgfstyle=straightmove, arrow=to, linewidth=0.07em, color=red!75!black, markmoves={e3-e4}]
}
\newcommand{\EbFen}{\fen{r3k2r/pbp2ppp/1pnqpn2/3p4/3P4/2PBPNP1/PP1N1PP1/R2QK2R w KQkq - 2 10}}
\newcommand{\EbMover}{White}
\newcommand{\EbPhase}{middlegame}
\newcommand{\EbStructure}{the London pyramid}
\newcommand{\EbBreak}{10.~e4}
\newcommand{\EbBreakSound}{sound}
\newcommand{\EbWpBreak}{55}
\newcommand{\EbBest}{10.~\symqueen{}a4}
\newcommand{\EbWpBest}{55}
\newcommand{\EbPlayed}{10.~\symknight{}e5}
\newcommand{\EbWpPlayed}{38}
\newcommand{\EbPlayedSound}{not sound}
\newcommand{\EbAimed}{\symbishop{}d3, \symknight{}d2 and \symqueen{}d1}
\newcommand{\EbAimedOk}{yes}
\newcommand{\EbUndeveloped}{}
\newcommand{\EbFiles}{d-file and e-file}
\newcommand{\EbRooks}{none}
\newcommand{\EbRooksOk}{no}
\newcommand{\EbForceUs}{18}
\newcommand{\EbForceThem}{12}
\newcommand{\EbForceOk}{yes}
\newcommand{\EbBoard}{%
\chessboard[setfen={r3k2r/pbp2ppp/1pnqpn2/3p4/3P4/2PBPNP1/PP1N1PP1/R2QK2R w KQkq - 2 10}, boardfontsize=14pt, pgfstyle=straightmove, arrow=to, linewidth=0.07em, color=blue!70!black, markmoves={f3-e5}, pgfstyle=straightmove, arrow=to, linewidth=0.07em, color=green!50!black, markmoves={e3-e4}]
}


\title{Does the checklist before a break tell when the break is sound?\\[0.3em]
\large The coaching method's five questions read as board code, on the \BreaksAll{} breaks of
his quiet positions and \StrongBreaks{} of strong players', refereed by Stockfish}
\author{Chess Coach research lab}
\date{}

\begin{document}
\maketitle

\begin{abstract}
A break the app teaches for his pawn structure, played merely because it can be, is sound about
one time in five, so what the app can teach is the moment. The coaching method the app follows
names the moment with five questions to ask before a break: are the pieces developed and aimed at
the lines it opens, are the rooks ready for the new file, does it create a passed pawn or expose
a king, is the force on that side equal or better, has it been calculated? Read with the app's
own board readings on the breaks of his quiet positions, all five pass for \AllFive{} of his
\BreaksAll{} breaks, so the rule tests the three questions about the pieces, among the breaks the
app's blunder check calls safe, comparing each break with the same break of the same structure at
other moments. By the rule fixed before any reading was joined to a score, there is no signal:
the lift is \LiftFirst{} and \LiftSecond{} points on the two halves of his games, with both
intervals across 0, and \StrongLift{} for strong players. The controls pass, and the blunder
check itself tells the moment: a safe break is sound about twice as often as one that is not
(\SafeSoundFirst\% and \SafeSoundSecond\% against \UnsafeSoundFirst\% and \UnsafeSoundSecond\%).
Even where a safe break has its pieces ready, it is sound \AgainstReadyFirst\% and
\AgainstReadySecond\% of the time, against his own move's \AgainstOwnFirst\% and
\AgainstOwnSecond\%. The app can keep the checklist as a way of thinking and put the blunder check
first, but the verdict on a break in a position stays the engine's.
\end{abstract}

\section{The question}

The study of the breaks the app teaches for each pawn structure found that a break played merely
because it can be is sound about one time in five in his quiet positions, no more often than a
random move, while his own move is sound more than half the time; the breaks he chose were as
sound as his other moves. What the app can teach is the moment. The coaching method the app
follows names the moment with a checklist to run before playing a break: pieces developed and
pointed at the lines that open, rooks ready for the new file, a passed pawn or an exposed king,
equal or better force on that sector, three or four moves calculated. It adds that rushing the
break before the pieces are ready is the most common structural error at his level.

Read as board code, does that checklist tell a sound break from an unsound one? If it does, the
review and the drills can teach the moment in the checklist's words: the pieces are not ready
yet, so the break opens lines for the other side. If it does not, the checklist read this way
cannot stand in for the engine's word on a position, and the app should not offer it as the test
of when to break.

\section{How it was measured}

\paragraph{The breaks.} Every break of every position the structure study scored: his quiet
middlegame and endgame positions with a break on, from his 10-minute games in the app's archive
at commit \texttt{\AppCommit}, in two halves by the games' end time (\PositionsFirst{} positions
in \GamesFirst{} games, and \PositionsSecond{} in \GamesSecond), and \StrongPositions{} positions
of strong players, one a game, from the Lichess games of September 2026 between two players rated
\Floor{} or more. Nothing is fitted, so the halves are two replications. A break is a legal pawn
move that the app's \texttt{playedLever} counts as a break of the structure the app names on the
board, for the side to move. Each break is a unit: his \BreaksFirst{} and \BreaksSecond{} breaks
and the strong players' \StrongBreaks. Its soundness is the structure study's: Stockfish at the
app's grading budget of \GradeNodes{} nodes, the hash cleared before each position, accepts it
when it gives up no more than \Window{} win\% against the best. Its group is its structure and the
break it plays, such as the London pyramid's c3-c4.

\paragraph{The questions, as board code.} Each question reads the position before the break,
through the app's knowledge service. The break moves a pawn to the break square; its targets are
the other side's pawns the pawn attacks from there; the lines it opens are the files of the
pawn's square, the break square and the targets; its sector is the break square's file and the
files on either side.
\begin{enumerate}[nosep]
\item \textbf{Pieces aimed}: no knight or bishop on its starting square, and at least \AimedMin{}
  of the mover's knights, bishops, rooks and queens aimed at the break: a rook or queen on one of
  the lines it opens, or a piece that sees the break square or a target.
\item \textbf{Rooks ready}: a rook on one of those lines, or with a legal move to one; a mover
  with no rook passes.
\item \textbf{A passer or an exposed king}: a new passed pawn for the mover, or a thinner shelter
  in front of the other king by the app's shelter scale, after the break or after the trade it
  offers (the break's pawn and a target off the board).
\item \textbf{Force on that side}: the mover's knights, bishops, rooks and queens on the sector's
  files are worth at least the other side's, by the app's piece values.
\item \textbf{Safe}: the app's blunder check calls the break safe: it allows no mate in one,
  leaves nothing newly hanging and deserts no defence. That is as much of three or four moves
  calculated as board code reads without an engine.
\end{enumerate}
A break whose pieces are ready passes questions 1, 2 and 4; a safe break whose pieces are not
ready is rushed.

\paragraph{The lift.} The lift of a reading is the share of sound breaks among those it passes
minus the share among those it fails, within each group, pooled over the groups with breaks on
both sides by the weight $n_1 n_0/(n_1+n_0)$. It compares each break with itself at other
moments, so it measures the moment, not which break is the sounder in general. Intervals come
from \Draws{} bootstrap draws with seed \Seed, resampling his games within a half, or the strong
players' positions: the 95\% percentile interval.

\section{Counted before the rule}

Before the rule was fixed, each question was read on every break and the passes counted, with no
score of any move read (Table~\ref{tab:coverage}). All five pass for \AllFive{} of his
\BreaksAll{} breaks and \AllFiveStrong{} of the strong players' \StrongBreaks: the question of a
passer or an exposed king passes for \PurposeShare\% of his breaks, and the checklist taken whole
cannot be tested. The rule tests the three piece questions, the ones the checklist calls the
common error, among the breaks the blunder check calls safe: \ReadyFirst{} of \SafeFirst{} and
\ReadySecond{} of \SafeSecond{} of his safe breaks have their pieces ready.

\CoverageTable

\section{The rule, fixed before any reading was joined to a score}

The controls come first, and if any fails there is no verdict. (1) The move played must be sound
more often than a random move that is not a break, by an interval above 0, on each half and for
the strong players. (2) A reading known to matter must show: the blunder check's own lift, over
all breaks, above 0 on both halves; without it, a lift the piece questions lack could be one this
design cannot see. (3) A reading drawn at random, passing as many safe breaks in each group as
the pieces' readiness does, must reach neither verdict below.

Then, on his safe breaks: the pieces tell the moment when the lift of their readiness lies above
0 on both halves; they mislead when it lies below 0 on both; otherwise there is no signal; and
there is not enough data when a half has fewer than \MinArm{} safe breaks with the pieces ready or
fewer than \MinArm{} rushed ones.

\section{The verdict: no signal}

The controls pass (Table~\ref{tab:controls}). His move is sound \ControlFirst{} and
\ControlSecond{} points more often than a random move, and the strong players' \StrongControl.
The blunder check tells the moment: \CheckFirst{} points (\CheckFirstCi) and \CheckSecond{}
(\CheckSecondCi), and \StrongCheck{} for the strong players. The random reading's lift is
\NullFirst{} (\NullFirstCi) and \NullSecond{} (\NullSecondCi): its intervals leave out 0 on both
halves, but on opposite sides, so it reaches neither verdict and the control passes by the rule.
Section~\ref{sec:after} checks those intervals on many random readings.

The pieces' readiness gives no signal: its lift is \LiftFirst{} points (\LiftFirstCi) and
\LiftSecond{} (\LiftSecondCi). The safe breaks with the pieces ready are sound \ReadySoundFirst\%
and \ReadySoundSecond\% of the time, the rushed ones \RushedSoundFirst\% and
\RushedSoundSecond\%. In the strong players' positions, reported and not ruled on, the lift is
\StrongLift{} (\StrongLiftCi).

\ControlTable

\section{Two positions}

Two of his positions show the checklist and the engine disagreeing (Figure~\ref{fig:examples}):
the first of the later half where a safe break with its pieces ready is unsound while his own move
is sound, and the first where a rushed safe break is sound while he played something else. They
illustrate; they test nothing.

\begin{figure}[ht]
\centering
\begin{subfigure}[t]{0.48\textwidth}
\centering
\EaBoard\\[0.3em]
{\footnotesize FEN \EaFen}
\caption{He is \EaMover, in the \EaPhase, in \EaStructure. The break \EaBreak{} opens the
\EaFiles. Aimed at it: \EaAimed\EaUndeveloped{} (\EaAimedOk); rooks ready: \EaRooks{}
(\EaRooksOk); force \EaForceUs{} against \EaForceThem{} (\EaForceOk); safe. Stockfish: \EaBest{}
is best at \EaWpBest\% to win; the break is \EaBreakSound{} (\EaWpBreak\%). He played \EaPlayed{}
(blue), \EaPlayedSound{} at \EaWpPlayed\%.}
\label{fig:ready}
\end{subfigure}\hfill
\begin{subfigure}[t]{0.48\textwidth}
\centering
\EbBoard\\[0.3em]
{\footnotesize FEN \EbFen}
\caption{He is \EbMover, in the \EbPhase, in \EbStructure. The break \EbBreak{} opens the
\EbFiles. Aimed at it: \EbAimed\EbUndeveloped{} (\EbAimedOk); rooks ready: \EbRooks{}
(\EbRooksOk); force \EbForceUs{} against \EbForceThem{} (\EbForceOk); safe. Stockfish: \EbBest{}
is best at \EbWpBest\% to win; the break is \EbBreakSound{} (\EbWpBreak\%). He played \EbPlayed{}
(blue), \EbPlayedSound{} at \EbWpPlayed\%.}
\label{fig:rushed}
\end{subfigure}
\caption{Two of his positions with a safe break. The break is green where Stockfish finds it
sound and red where it does not; win\% is his after each move. Stockfish at \GradeNodes{} nodes,
rerun for this figure.}
\label{fig:examples}
\end{figure}

\section{What the breaks describe}

\paragraph{Each question alone.} Over all the breaks, safe or not, the blunder check has the
largest lift, \QSafe{} points with his halves together (Table~\ref{tab:questions}), and for the
strong players it is the only question whose interval lies above 0. For him, the pieces aimed
(\QAimed{} points) and the force on that side (\QForce) also lie above 0 over all breaks, safe or
not; the rooks ready and the question of a passer or an exposed king do not.

\QuestionTable

\paragraph{The more questions pass.} Over all his breaks, the share sound rises with the number
of questions passed: \CountOne\% with one, \CountTwo\% with two, \CountThree\% with three and
\CountFour\% with four (Table~\ref{tab:count}). That count mixes the groups and the blunder
check, and four questions passing still leave the break sound less often than his own move.

\CountTable

\paragraph{By group.} In each of the three groups with enough safe breaks on both sides, the
London pyramid's two breaks and the advance chain's \ldots f7-f6, the lift holds 0
(Table~\ref{tab:groups}); the other \GroupsOther{} groups have too few breaks on one side.

\GroupTable

\paragraph{His choice, and the ready break against his move.} He played a safe break with the
pieces ready \PlayedReady\% of the time, a rushed one \PlayedRushed\% and an unsafe one
\PlayedUnsafe\% (Table~\ref{tab:choice}). Within the groups, the share of ready breaks he
played minus the share of rushed ones is \ChoiceLift{} points (\ChoiceLiftCi): he does not wait
for the pieces, but he does avoid the unsafe breaks, as the strong players do. Where a safe break has its pieces ready, in \AgainstFirst{} and
\AgainstSecond{} of his positions, those breaks are sound \AgainstReadyFirst\% and
\AgainstReadySecond\% of the time and his own moves there \AgainstOwnFirst\% and
\AgainstOwnSecond\%: the ready break trails his move by \AgainstGainFirst{} points
(\AgainstGainFirstCi) and \AgainstGainSecond{} (\AgainstGainSecondCi).

\ChoiceTable

\section{Added after the run}
\label{sec:after}

\paragraph{The intervals.} The random reading's intervals left out 0 on both halves, which a
95\% interval should do for about one random reading in twenty on each. Over
\CalPermutations{} random readings drawn the same way, the lift spreads by \CalSpreadFirst{} and
\CalSpreadSecond{} points, and the bootstrap's own spread for \CalReadings{} of them is
\CalBootFirst{} and \CalBootSecond; their intervals leave out 0 for \CalExcludesFirst{} and
\CalExcludesSecond{} of the \CalReadings. The intervals hold their width; the reading the rule
drew lies \NullSdFirst{} and \NullSdSecond{} spreads from 0, an unlucky draw. The pieces'
readiness lies no farther from 0 than random readings often do: \CalPFirst{} and \CalPSecond{} of
them lie as far (Table~\ref{tab:safequestions}).

\paragraph{Each piece question alone, among the safe breaks.} With his halves together, the pieces
aimed give \SAimed{} points (\SAimedCi) and the force on that side \SForce{} (\SForceCi), while
the rooks give \SRooks{} (\SRooksCi); for the strong players none lies above 0 (\SAimedStrong,
\SRooksStrong{} and \SForceStrong). Asking for all three loses what the aimed pieces and the force
carry: his safe breaks passing two of the piece questions are sound \SafeTwo\% of the time and
those passing all three \SafeThree\% (Table~\ref{tab:count}). These readings were chosen after
the run, among several, so they do not overturn the verdict; a rule built on the aimed pieces or
the force needs positions this study did not see.

\SafeQuestionTable

\section{What the app can take}

\begin{itemize}[nosep]
\item \textbf{The blunder check first.} Of the checklist's questions, the one that tells the
  moment is the calculation, as the app's blunder check reads it: a safe break is sound about
  twice as often as one that is not, for him and for strong players. When the review, a drill or
  the coach's reading of a position names a break, the first thing it checks is that the break
  leaves nothing loose.
\item \textbf{Not the piece questions as the verdict.} With the break safe, whether the pieces are
  aimed, the rooks ready and the force there tells nothing about whether the break is sound, for
  him or for strong players. The app should not call a break ready or rushed by these readings
  as a judgment of the move. The questions can stay as a way of thinking about a structure; the
  verdict on a break in a position is Stockfish's.
\item \textbf{A ready break is still no suggestion.} Where a safe break has its pieces ready, it
  is sound about three times in ten, his own move more than five times in ten. The app should
  never offer a break as the move on the checklist's word.
\end{itemize}

\section{Caveats}

\begin{itemize}[nosep]
\item \textbf{A reading of the words.} ``Pointed at the lines that open'', ``ready'' and
  ``sector'' are read here one way each, fixed before any score was joined; a coach might read
  them otherwise, which would be another study.
\item \textbf{The checklist taken whole.} All five questions almost never pass, so only the piece
  questions and each question apart were tested.
\item \textbf{The breaks are the app's.} Only the breaks the structure catalogue names are tested,
  at the squares the app reads.
\item \textbf{Sound is a window.} Several moves are often within \Window{} win\% of the best; a
  break that is not the best can be sound.
\item \textbf{Many breaks a game for him.} A game where a break stays on counts many times;
  intervals resample whole games for that reason.
\item \textbf{The same positions.} The positions, scores and referee are the structure study's;
  only the two positions shown were searched again.
\item \textbf{Added after the run.} The check of the intervals and the piece questions alone among
  the safe breaks were chosen once the rule had been applied. They describe and decide nothing.
\end{itemize}

\end{document}
